\section{Related Work} \label{sec:related_work}

Interactive theorem provers such as 
 Lean 4~\cite{Moura2021}, 
 Rocq (formerly Coq)~\cite{rocq-reference-manual}
and 
 Isabelle/HOL~\cite{PaulsonNipkowWenzel2019} allow definitions 
 and proofs to be stated in a formal
language and machine-checked, rather than only reviewed by hand. 
They have
been used to formalize substantial bodies of mathematics and, increasingly,
the semantics of programming languages and logics used in practice~\cite{HuangEtAl2026}, 
as illustrated by verified systems such as 
the CompCert compiler~\cite{Leroy2009CompCert,Leroy2012MechanizedSemantics}, 
the seL4 microkernel~\cite{KleinEtAl2009}, 
and by more recent industrial applications such as AWS's Cedar authorization language, 
formalized in Lean~\cite{CutlerEtAl2024}.

The application of these techniques to Knowledge Graphs and the Semantic Web is still less mature than in other branches of mathematics.
Some earlier work used Isabelle/HOL to formalize OWL semantics
 in \cite{TurnerCarrol07}, and
 \cite{tantow_et_al:LIPIcs.ITP.2025.36} verifies Datalog using Lean.

ShEx was initially introduced in~\cite{PGS14} in an informal way based on an implementation. 
Later,~\cite{SBG15} established the complexity of Shape Expressions without negation 
and presented a formal semantics for them based on regular bag expressions. An implementation based on regular expression derivatives was proposed in~\cite{labragayo2015towards}.
 The formal semantics of the ShEx language including negation and recursion was presented in~\cite{BGP17} 
  and published as a W3C Community Group Report~\cite{shex_standard}. 
 More recently,~\cite{ShExInheritance} presented an extension of ShEx with inheritance,
  which is currently developed by the IEEE P3330 Shape Expression Schemas Working Group.
  Subsequent work has compared ShEx with other validation approaches,
 such as SHACL~\cite{SHACL} and PGSchema~\cite{ABDF23},
 and formalized the common foundations of the three languages:
  ~\cite{www25} compares ShEx, SHACL and PGSchema without considering recursion,
  and~\cite{recursive_shapes} compares ShEx and SHACL with recursion. None of these works employed an automatic theorem prover to verify
  the correctness of their semantics and algorithms.
